\documentclass[10pt]{article}
\usepackage[margin=20mm]{geometry}
\usepackage{longtable,booktabs,array}
\usepackage{graphicx,url}
\usepackage{placeins}
\usepackage{amsmath,amssymb}
\begin{document}
\section*{Supplement C: Personnel profiles and comparison designs}
\subsection*{C1. Person-by-difficulty results}
For each person and difficulty, the cell is number of images / image-mean Q / mean relative rank within the image. A larger rank means higher relative Q. Dashes indicate unavailable values. These are descriptive primary-target Manual records; historical profile construction and the same-room predictions use their own support.
\begin{longtable}{llll}
\toprule Annotator & Easy: N / Q / rank & Moderate: N / Q / rank & Difficult: N / Q / rank \\
\midrule\endfirsthead\toprule Annotator & Easy: N / Q / rank & Moderate: N / Q / rank & Difficult: N / Q / rank \\
\midrule\endhead
A1 & 44 / 68.608 / 0.336 & 9 / 59.374 / 0.436 & 16 / 48.090 / 0.435 \\
A2 & 34 / 79.108 / 0.583 & 9 / 84.622 / 0.894 & 15 / 80.054 / 0.866 \\
A3 & 44 / 73.648 / 0.447 & 10 / 58.964 / 0.433 & 16 / 51.298 / 0.641 \\
A4 & 45 / 78.187 / 0.577 & 10 / 71.358 / 0.697 & 19 / 60.233 / 0.693 \\
A5 & 41 / 75.511 / 0.505 & 12 / 60.596 / 0.513 & 16 / 49.310 / 0.542 \\
A6 & 27 / 72.485 / 0.493 & 6 / 63.396 / 0.499 & 14 / 48.266 / 0.566 \\
A7 & 44 / 59.648 / 0.173 & 8 / 60.437 / 0.495 & 14 / 39.769 / 0.343 \\
A8 & 35 / 71.732 / 0.452 & 9 / 57.257 / 0.509 & 18 / 50.025 / 0.574 \\
A9 & 21 / 73.043 / 0.483 & 5 / 55.203 / 0.521 & 11 / 44.398 / 0.521 \\
A10 & 46 / 79.222 / 0.566 & 11 / 72.963 / 0.633 & 16 / 60.502 / 0.695 \\
A11 & 46 / 73.865 / 0.446 & 7 / 60.876 / 0.480 & 19 / 63.060 / 0.824 \\
A12 & 20 / 80.491 / 0.635 & 5 / 83.815 / 0.917 & 11 / 59.461 / 0.615 \\
A14 & 12 / 77.446 / 0.541 & 2 / 46.817 / 0.184 & 6 / 54.171 / 0.704 \\
A15 & 2 / 74.496 / 0.314 & -- & -- \\
A17 & 23 / 86.152 / 0.806 & 3 / 74.981 / 0.780 & 9 / 72.165 / 0.816 \\
A18 & 47 / 71.879 / 0.432 & 10 / 54.242 / 0.376 & 15 / 47.106 / 0.475 \\
A19 & 43 / 78.774 / 0.586 & 8 / 54.146 / 0.374 & 13 / 41.168 / 0.428 \\
A20 & 48 / 77.535 / 0.523 & 9 / 60.038 / 0.471 & 19 / 48.385 / 0.479 \\
A21 & 55 / 72.895 / 0.458 & 10 / 65.684 / 0.608 & 16 / 41.927 / 0.458 \\
A22 & 50 / 67.729 / 0.380 & 10 / 56.921 / 0.452 & 15 / 46.204 / 0.472 \\
A23 & 47 / 79.078 / 0.589 & 11 / 67.775 / 0.607 & 18 / 59.241 / 0.737 \\
A24 & 47 / 70.758 / 0.398 & 8 / 56.128 / 0.482 & 20 / 39.536 / 0.347 \\
A25 & 50 / 80.740 / 0.593 & 9 / 55.117 / 0.400 & 15 / 45.698 / 0.454 \\
A26 & 44 / 83.431 / 0.682 & 9 / 69.788 / 0.577 & 15 / 48.704 / 0.496 \\
A27 & 39 / 68.003 / 0.302 & 7 / 38.185 / 0.182 & 10 / 41.174 / 0.369 \\
A28 & 91 / 73.479 / 0.422 & 28 / 57.761 / 0.466 & 20 / 43.921 / 0.385 \\
A29 & 91 / 71.946 / 0.390 & 28 / 56.543 / 0.386 & 19 / 47.023 / 0.455 \\
A30 & 87 / 91.038 / 0.783 & 25 / 62.480 / 0.535 & 21 / 68.581 / 0.747 \\
A31 & 88 / 72.038 / 0.386 & 26 / 55.658 / 0.406 & 20 / 43.369 / 0.354 \\
A32 & 89 / 76.792 / 0.476 & 27 / 68.907 / 0.596 & 21 / 48.774 / 0.509 \\
A33 & 90 / 71.316 / 0.351 & 30 / 53.453 / 0.372 & 19 / 45.935 / 0.452 \\
A34 & 91 / 70.420 / 0.360 & 27 / 53.882 / 0.372 & 19 / 42.974 / 0.388 \\
A35 & 90 / 73.377 / 0.414 & 27 / 64.279 / 0.518 & 21 / 43.677 / 0.351 \\
A36 & 88 / 76.754 / 0.482 & 28 / 67.296 / 0.591 & 21 / 40.146 / 0.314 \\
A37 & 87 / 92.386 / 0.819 & 26 / 72.056 / 0.714 & 21 / 58.001 / 0.592 \\
A38 & 93 / 75.924 / 0.469 & 28 / 65.312 / 0.576 & 22 / 50.460 / 0.534 \\
A39 & 93 / 69.970 / 0.336 & 28 / 58.390 / 0.469 & 19 / 46.983 / 0.473 \\
A40 & 86 / 92.660 / 0.828 & 25 / 68.145 / 0.579 & 20 / 66.457 / 0.812 \\
A41 & 95 / 71.744 / 0.403 & 27 / 54.366 / 0.410 & 22 / 45.524 / 0.436 \\
A42 & 94 / 78.580 / 0.523 & 28 / 61.210 / 0.504 & 22 / 40.871 / 0.310 \\
A43 & 87 / 93.871 / 0.852 & 27 / 71.636 / 0.655 & 22 / 61.157 / 0.667 \\
A44 & 94 / 76.137 / 0.492 & 28 / 60.665 / 0.478 & 20 / 46.546 / 0.460 \\
A45 & 93 / 72.174 / 0.412 & 27 / 56.678 / 0.421 & 19 / 42.190 / 0.366 \\
A46 & 93 / 79.039 / 0.512 & 28 / 60.370 / 0.414 & 20 / 41.665 / 0.408 \\
A47 & 94 / 71.388 / 0.375 & 29 / 52.207 / 0.365 & 22 / 44.175 / 0.402 \\
A48 & 94 / 74.922 / 0.442 & 28 / 60.826 / 0.509 & 21 / 41.475 / 0.330 \\
A49 & 94 / 69.966 / 0.373 & 27 / 50.477 / 0.328 & 19 / 44.359 / 0.435 \\
A50 & 89 / 92.232 / 0.834 & 26 / 67.604 / 0.611 & 22 / 63.241 / 0.725 \\
A51 & 94 / 77.473 / 0.494 & 28 / 62.769 / 0.494 & 24 / 46.318 / 0.436 \\
A52 & 95 / 73.065 / 0.408 & 29 / 49.498 / 0.314 & 23 / 38.766 / 0.393 \\
A53 & 95 / 79.230 / 0.531 & 30 / 61.352 / 0.519 & 19 / 44.528 / 0.406 \\
A54 & 93 / 89.109 / 0.780 & 28 / 73.939 / 0.675 & 21 / 67.382 / 0.787 \\
A55 & 97 / 74.701 / 0.451 & 30 / 64.548 / 0.575 & 22 / 38.963 / 0.361 \\
A56 & 100 / 71.618 / 0.383 & 30 / 54.035 / 0.404 & 22 / 43.471 / 0.377 \\
A57 & 100 / 71.689 / 0.376 & 30 / 57.006 / 0.449 & 22 / 41.974 / 0.359 \\
A58 & 99 / 91.628 / 0.804 & 30 / 71.334 / 0.630 & 22 / 63.183 / 0.728 \\
A59 & 100 / 78.011 / 0.505 & 31 / 64.492 / 0.508 & 25 / 44.475 / 0.439 \\
A60 & 102 / 72.485 / 0.400 & 32 / 58.019 / 0.470 & 26 / 42.354 / 0.430 \\
A61 & 103 / 71.957 / 0.395 & 31 / 54.177 / 0.389 & 24 / 46.705 / 0.438 \\
A62 & 103 / 76.474 / 0.473 & 32 / 69.563 / 0.655 & 22 / 46.100 / 0.394 \\
\bottomrule\end{longtable}
\subsection*{C2. Structural selection versus historical-quality selection}
The matched selection comparison covers 208 events and 670 distinct teams. At 14 people with seven replaced, structural selection minus historical-quality selection averaged $-$0.034 Q points. At 20 people with seven replaced, the difference was $-$0.482; the building-cluster interval was [$-$1.220, $-$0.001]. This comparison is distinct from the quality-and-time policy in the main text. Target-team identities, output status, and shared support follow the recorded comparison.
\subsection*{C3. Matched replacements and historical profiles}
\begin{center}\small
\begin{tabular}{lrrrr}
\toprule Design / profile & Team $n$ & Replaced $m$ & Images & Mean C--B Q \\
\midrule
Joint profile K4, recorded 12/6 & 12 & 6 & 56 & +8.051 \\
Consensus-response K2 & 12 & 4 & 83 & see note \\
Consensus-response K2 & 20 & 4 & 82 & see note \\
Consensus-response K2 & 26 & 3 & 112 & see note \\
More-simplifying to turn-retaining & 8 & 4 & 24 & +10.095 \\
Reverse direction & 8 & 4 & 4 & $-$5.404 \\
\bottomrule
\end{tabular}\end{center}
The joint-profile row retains its recorded class labels and historical design. The consensus-response rows use the R1/R2 display aliases defined in the main text and retain the recorded class memberships. They use recorded matched comparisons; within- and between-profile absolute differences are described in the main text for 12/4, not substituted with a directional difference. The two eight-person directions have different image supports.

\subsection*{C4. Large replacement blocks}
The table summarizes image-level results for the stated profile definitions and class labels. $m$ is the number replaced and $n$ is team size; image counts vary with capacity. Values are image-equal mean C--B Q among the recorded image-level results.
\begin{longtable}{llrrrrr}
\toprule Profile method & Direction & $n$ & $m$ & Images & Buildings & Mean C--B \\
\midrule\endfirsthead\toprule Profile method & Direction & $n$ & $m$ & Images & Buildings & Mean C--B \\
\midrule\endhead
Target structural K3 & C2 to C1 & 11 & 9 & 8 & 4 & -6.771 \\
Target structural K3 & C2 to C1 & 12 & 10 & 1 & 1 & +1.596 \\
Target structural K3 & C2 to C1 & 14 & 9 & 1 & 1 & -1.617 \\
Target structural K3 & C2 to C1 & 14 & 10 & 1 & 1 & +2.277 \\
Target structural K3 & C2 to C1 & 15 & 9 & 1 & 1 & -28.367 \\
Target structural K3 & C2 to C1 & 16 & 9 & 1 & 1 & +1.408 \\
Target structural K3 & C2 to C1 & 17 & 9 & 5 & 3 & -8.428 \\
RAW3 structural K3 & C2 to C1 & 11 & 9 & 15 & 8 & -4.533 \\
RAW3 structural K3 & C2 to C1 & 13 & 9 & 1 & 1 & -0.406 \\
RAW3 structural K3 & C2 to C1 & 14 & 9 & 2 & 2 & +0.886 \\
RAW3 structural K3 & C2 to C1 & 15 & 9 & 1 & 1 & -1.612 \\
RAW3 structural K3 & C2 to C1 & 16 & 9 & 2 & 2 & -0.156 \\
RAW3 structural K3 & C2 to C1 & 17 & 9 & 9 & 5 & -4.334 \\
RAW3 structural K3 & C2 to C3 & 11 & 9 & 8 & 6 & +9.015 \\
RAW3 structural K3 & C2 to C3 & 12 & 10 & 1 & 1 & +8.330 \\
RAW3 structural K3 & C2 to C3 & 15 & 10 & 1 & 1 & +7.114 \\
RAW3 structural K3 & C2 to C3 & 17 & 9 & 8 & 6 & +6.054 \\
\bottomrule\end{longtable}
Class labels retain the definitions in each listed design. Some conditions contain only one image, so the table reports both capacity and the corresponding quality change.
\clearpage
\subsection*{C5. Floor-layout variation across teams}
The same 171-image panel and 256 entry orders per image underlie these curves. Each point summarizes pairwise floor differences between teams at the same group size, with equal image weighting. This geometric quantity complements the score variation in the main text.
\begin{figure}[ht]
\centering
\includegraphics[width=.95\linewidth]{supplement_figures/floor_diversity.pdf}
\caption{Floor-layout variation for all images and the three difficulty groups. Doorway and OOS analyses are separate.}
\end{figure}

% Insert before \end{document} in supplement_personnel.tex.
% Uses only its existing longtable, booktabs and array packages.
% MR source IDs resolve to the accompanying source_manifest.json and SOURCE_MAP.txt.
\clearpage
\subsection*{C6. Reproducible personnel-profile specifications}
The specifications below define the implementations used for the reported
structural-expression, consensus-response and joint-profile analyses
[MR01--MR03]. For an observation in building
$b$, labels are joined by person and \texttt{heldout\_building}$=b$.
The \texttt{ALL} fit is descriptive only. Class numbers are assigned by the
ordered training-fold class centres, not by matching a test fold to the
\texttt{ALL} solution. Other replacement designs retain the labels specified
for their own comparisons.

Let $S$ denote four raw-response fields:
\begin{sloppypar}
\texttt{turn\_n}, \texttt{short\_turn\_edges},
\texttt{near\_\allowbreak straight\_\allowbreak nodes\_\allowbreak 5deg},
and \texttt{unmatched\_\allowbreak ref\_\allowbreak 9}.
They count floor turns $>5^\circ$, edges shorter than 4\% of the perimeter
with both endpoint turns at least $15^\circ$, turns $\leq5^\circ$, and the
unmatched reference fraction after one-to-one spherical floor-point matching
with a $9^\circ$ acceptance threshold, respectively.
Let $B$ denote \texttt{matched\_\allowbreak top\_\allowbreak bias\_\allowbreak deg},
\texttt{matched\_\allowbreak bottom\_\allowbreak bias\_\allowbreak deg}, \texttt{log\_area\_ratio}, and
\texttt{log\_height}. The signed endpoint deviations use one-to-one periodic
azimuth matching with a $9^\circ$ threshold; the remaining fields are log
submitted/reference floor area and log reconstructed mean height under
camera-height normalization. These are not all absolute-error measures.
\end{sloppypar}

Raw profile features are centred on the median for image, legal reference
and condition. Log time additionally separates origin and time source and
uses positive finite recorded durations only. Each metric is averaged over
valid records within person--target, then equally over targets within image,
images within person--building, and buildings within person. Missing
outcomes are not zero-imputed [MR02].

\begingroup\small
\begin{longtable}{>{\raggedright\arraybackslash}p{.19\linewidth}>{\raggedright\arraybackslash}p{.49\linewidth}>{\raggedright\arraybackslash}p{.23\linewidth}}
\toprule
Profile & Exact feature blocks & Eligibility and class order \\
\midrule\endfirsthead
\toprule Profile & Exact feature blocks & Eligibility and class order \\
\midrule\endhead
Target structural
(\texttt{TARGET\_DSTRUCT})
&
12 features, two domains: overall means of $S$ (4); difficult means and
difficult-minus-easy means of $S$ (8).
One Ward tree has the defined $K=2,3,4$ cuts.
&
Overall $\geq12$ images/3 buildings; easy $\geq5$/2; difficult $\geq4$/2.
Complete required features.
Order by difficult turn count, difficult short-edge count, overall turn count.
\\[3pt]
Consensus response
(\texttt{RFUSE\_GEOM})
&
21 regularized member-effect features: overall, difficult and
difficult-minus-easy coefficients for
\texttt{floor\_turn5\_n},
\texttt{floor\_short\_edge\_03\_n},
\texttt{floor\_redundant5\_n} (structure domain, 9);
\texttt{bottom\_y\_mean\_px}, \texttt{top\_y\_mean\_px},
\texttt{log\_floor\_area}, \texttt{log\_height\_mean}
(boundary/scale domain, 12).
The paper uses the defined $K=2$ cut.
&
Overall $\geq8$ images/3 buildings; easy $\geq5$/2; difficult $\geq3$/2.
Order by difficult fused turn response, then overall bottom-coordinate response.
Native fused-output short edges use the length threshold defined below,
distinct from $S$.
\\[3pt]
Joint profile
(\texttt{OPPJOINT})
&
13 features, four domains: opportunity slopes of $S$ (4), opportunity slopes
of $B$ (4), overall mean/p20 and opportunity slope of peer-centred Q (3),
overall mean and opportunity slope of peer-centred log time (2).
The reported $K=3,4$ are existing cuts.
&
Overall $\geq12$ images/3 buildings; opportunity variance $>10^{-5}$ and
at least two levels; positive-time support $\geq12$/3 with a defined time slope.
Complete required features.
Order by overall Q residual, challenge-turn residual, overall log-time residual.
\\
\bottomrule
\end{longtable}
\endgroup

For valid positive-area fused footprints, the short-edge count is the number
of cyclic edges with Euclidean length strictly below 0.3 camera-height units.
It applies neither a perimeter-normalized threshold nor an endpoint-turn
threshold [MR12].

For the joint profile, the opportunity inputs are
$\log(1+\max(v-4,0))$, $\log(1+H)$ and $\log(1+s)$, where
$v$, $H$ and $s$ are the archived \texttt{ref\_turn5}, \texttt{ref\_H} and
\texttt{ref\_short\_edges} fields.
Training unique image--reference targets receive equal
target-within-image, image-within-building and building weights.
Each axis is centred on its weighted median, divided by its weighted
p85--p15 range (fallback: weighted standard deviation, then 1), clipped to
$[-3,3]$, and the three axes are summed and divided by $\sqrt{3}$.
A person's response slope is weighted covariance of the peer residual and
opportunity divided by weighted opportunity variance. The Q p20 is a weighted
empirical quantile of the same peer residuals, not a raw-Q P10 [MR02].

For each eligible training-fold profile domain, columns are median-centred,
scaled by p85--p15 (fallback: population standard deviation, then 1), and
clipped to $[-3,3]$. Full-SVD PCA retains at most two axes per raw-response
domain and two per consensus-response domain. Each domain is divided by
the training root-mean-squared PC-vector length (floor $10^{-9}$) and by
the square root of the number of domains. Ward linkage uses Euclidean
distances and optimal leaf ordering. The target building is excluded before
history, eligibility, scaling, PCA and clustering. Insufficient histories
remain \texttt{unknown}; they are not assigned to the nearest class.

\paragraph{Regularized consensus-response effects.}
The \texttt{RFUSE\_GEOM} implementation uses ordinary Manual natural
orders, sizes $k\in\{10,12,14,16,20,24\}$, and 24 natural orders per pool.
Repeated deterministic team endpoints are collapsed with their path
occurrence counts retained. A pool--size group requires at least four
distinct eligible teams. The seven geometry outputs and the shared
quality-output eligibility fields must be available; thus Q is not a
GEOM clustering feature, but Q availability affects its training support.
For team $t$ and person $w$, $X_{tw}=1/k_t$ for a member and zero otherwise.
Within pool--size groups, both $X$ and the output matrix $Y$ are centred
using path-occurrence weights. Each building has total weight one, with
equal images within building and equal target--size groups within image;
path occurrences determine weights within each group. The separate
overall/easy/difficult fits are
\[
 \widehat{\mathbf B}
 = (X^\mathsf{T}WX+\lambda I)^{-1}X^\mathsf{T}WY,\qquad
 \lambda=\max\{ \alpha\,\mathrm{tr}(X^\mathsf{T}WX)/P,\;10^{-12}\},
\]
where $P$ is the number of membership columns.
Nested leave-one-building-out selection considers
$\alpha\in\{0.1,1,10\}$ using training-variance-normalized response MSE,
balanced over validation buildings, outputs, strata and domains.
The selection table chooses $\alpha=10$ for the descriptive fit
and all 22 outer folds [MR03]. This is the dimensionless multiplier, not
a fixed numerical penalty of 10. The target-equal
\texttt{TARGET\_RFUSE} and Q-augmented \texttt{RFUSE\_GEOMQ} candidates are
different models and are not substituted here.

\subsection*{C7. Prediction models, preprocessing and test support}
The following rows identify the exact implementations behind the reported
comparisons. A shared label such as \emph{history} does not imply identical
aggregation across these analyses [MR04--MR07].

\begingroup\small
\begin{longtable}{>{\raggedright\arraybackslash}p{.19\linewidth}>{\raggedright\arraybackslash}p{.49\linewidth}>{\raggedright\arraybackslash}p{.23\linewidth}}
\toprule
Comparison & Features, estimator and preprocessing & Split, support and estimand \\
\midrule\endfirsthead
\toprule Comparison & Features, estimator and preprocessing & Split, support and estimand \\
\midrule\endhead
Difficulty versus computable geometry
&
Weighted linear ridge, penalty 1 on every coefficient except the intercept.
Controls: training-level person and record-source indicators; medium/hard
difficulty indicators. Added numeric features:
\texttt{ref\_turn5}, \texttt{ref\_H}, \texttt{ref\_short\_edges},
\texttt{ref\_min\_horizontal\_deg}, short-edge presence, and minimum
short-edge span conditional on presence (zero for structural absence).
Numeric weighted mean/SD scaling is training-only; no tuning.
The binary simplification prediction is clipped linear regression,
not logistic regression.
&
22 building-held-out folds; 6,512 records, 206 images, 60 people.
Training images have equal total weight; test errors are averaged within
image then building, with buildings equal.
MSE for turn deficit and simplification; MAE for Q.
Missing outcomes are excluded per outcome; nonfinite geometry is excluded.
Unseen people have zero person indicators.
\\[3pt]
Same-room image-mean Q:
9.496 to 6.464
&
Direct source-image mean Q substitutes for the target-image mean.
The baseline is the mean of other-building image means.
No fitted regression coefficient, scaling, imputation or hyperparameter.
&
128 directly adopted, non-pending pairs, both directions:
256 predictions, 65 images, 13 buildings.
Average absolute errors over source directions per target,
then images within building and buildings equally.
\\[3pt]
Same-room individual Q:
15.645 to 13.385
&
Baseline: the person's record-mean Q outside the target building.
Let $\theta_w$ be the record-mean outside-building residual from
leave-one-person-out task-peer Q.
Prediction: source-image mean Q of other people plus $\theta_w$,
clipped to $[0,100]$.
This is \texttt{public\_external} with a single direct source,
not the half-own-residual variant.
&
1,793 target answers, 65 images, 56 people, 13 buildings;
5,198 source-specific prediction rows.
Shared valid support across the saved predictors; missingness is not zero.
Source-specific errors are averaged per target answer before
image/building aggregation.
\\[3pt]
Same-room structure:
short edges 0.526 to 0.472
&
For each target answer, average the person's feature over directly
neighbouring source views on which that person answered.
The comparator averages other people's features on the same source views.
No fitted parameters. This averages source features before evaluating
the target error, unlike the single-direct Q comparison.
&
Same 1,793/65/56/13 support, requiring outside-building personal history.
The signed floor-offset comparison is in ERP pixels and uses spherical
bottom-point matching; it is not the azimuth-paired degree feature in $B$.
\\[3pt]
Recorded purpose added to geometry:
5.158 versus 5.379
&
Five numeric features: reference turn count, log(1+\texttt{ref\_H}),
short-edge/turn fraction, log minimum horizontal span and log
90th-percentile floor radius (last two floored at $10^{-5}$ before log).
Difficulty is categorical; the augmented model adds purpose.
Training-median imputation with missing indicators, standard scaling,
one-hot categories with unseen levels ignored.
Fixed candidate portfolio detailed below.
&
Outer LOBO on 206 ordinary images/22 buildings; reported repeated-purpose
subset: 190 images/21 buildings.
Inner building-grouped five-fold model selection.
Q predictions clipped to $[0,100]$; reported MAE is image-equal.
Geometry is reference-derived, not a no-reference RGB predictor.
\\
\bottomrule
\end{longtable}
\endgroup

The recorded-purpose portfolio contains linear ridge
$\alpha\in\{0.1,1,10,100\}$, quadratic-numeric ridge
$\alpha\in\{10,100\}$, and histogram gradient boosting
(120 iterations, 7 leaf nodes, minimum 15 samples per leaf,
learning rate 0.05, L2 penalty 10, seed 2026102703).
Quadratic expansion follows imputation/scaling and has degree two with no
bias column. The inner criterion is the mean of held-out-fold MAEs.
Every preprocessing step is re-fitted inside each inner training fold.
For Q, the geometry-only model selected quadratic ridge $\alpha=10$ in
16 outer folds and boosting in 6; geometry plus purpose selected quadratic
ridge $\alpha=10$ in 18, boosting in 3, and linear ridge $\alpha=1$ in 1.
The archived 44-row model table identifies every fold [MR07].

\subsection*{C8. Continuous Q and Q/T selection}
The fixed candidate roster is shared by the policies and does not expand
when the historical selection criteria change.
All target-building records are excluded from historical estimates.
For each person, $q_w$ is the record-mean Q residual relative to other people
on the same image, condition, legal target and reference version.
The cost $c_w$ is the exponential of the record-mean log-time residual,
additionally separating origin and time source. Q and time histories
each require at least eight images and three buildings.
The main panel uses the compatible complete-time history subset, without
mixing incompatible time sources [MR08].

Q-only selects the $n$ largest $q_w$, breaking ties by person identifier.
Q/T starts from that team and repeatedly chooses the best strictly
cost-reducing one-person exchange subject to a fixed historical mean-Q
floor of the initial team's mean minus 2 points. Feasible exchanges are
ordered by total $c_w$, then descending mean $q_w$, then sorted person IDs;
cost must decrease by more than $10^{-9}$. It stops when no such exchange
exists. This is a deterministic local-search rule, not a globally optimal
solution, and target Q does not select the team. Each
$n\in\{3,4,5,6,7,9,13\}$ is optimized separately, with more than $n$
eligible candidates; teams across sizes need not be nested.

The primary common-support identity is
(source--target pair, source pool, target pool, domain).
Both policies must have been executed at every listed size, with exactly
the same candidate roster throughout. Only ordinary Manual, comparable,
directly adopted relations without pending or substantive scope conflicts
are included. This yields 126 contexts, 46 target images and 9 buildings.
Let $\mathcal A$ average contexts within target image, images within
building, and buildings equally. The two primary estimands are
\[
 \mathcal A(Q_{\rm QT}-Q_{\rm Q}),\qquad
 \exp\{\mathcal A[\log(T_{\rm QT}/T_{\rm Q})]\}.
\]
Thus the labor ratio is not a ratio of arithmetic mean labor.
At $n=5$, these are $0.21836$ and $0.81343$, with empirical
building-cluster percentile intervals $[-0.53836,0.86896]$ and
$[0.76045,0.88192]$. The archive uses 20,000 building-bootstrap
resamples, seed 20261030, synchronized across sizes [MR09].
These quantify empirical cluster sensitivity, not a guarantee of
population-sampling coverage.

The $n=5$ versus $n=13$ comparison additionally requires both policies to
be scoreable at every size: 124 contexts, the same 46 images and 9 buildings.
Its labor-growth estimand is
$\exp\{\mathcal A[\log(T_{13}/T_5)]\}$, separately for each policy
(2.33142 and 2.17016). It compares independently selected teams,
not the causal effect of adding eight people to a fixed team.
Unscoreable Q is never replaced by zero.

\subsection*{C9. Source-history stopping and acquisition cost}
The reported 86.054 versus 85.189 Q comparison is the archived
\texttt{maximum\_history}, \texttt{whole\_group},
\texttt{all\_attempts}, non-pending-direct-relations panel.
For a room group containing $v$ views, it uses the maximum available
source-view history $m=v-1$, not one source view [MR10].
For each shared member-order replay $r$, let
\[
 s_r(k)=\frac{1}{m}\sum_{h\in H} Q_{hr}(k),\qquad
 M_r=\max_{j\geq3:\,s_r(j)\ {\rm finite}}s_r(j).
\]
\[
 k_{\rm req,r}=
 \min\{k\geq3:s_r(k)\ {\rm finite},\ s_r(k)\geq M_r-2\}.
\]
All source values at a prefix must be available; the implementation uses
ordinary means, no interpolation, and a numerical comparison tolerance
of $10^{-10}$. If no valid prefix exists, no choice is produced.
The target then advances from $k_{\rm req,r}$ to its first complete
geometric output: status \texttt{ok}, at least four nodes, and an available
connection. A target with no subsequent complete output remains unavailable.
Target reference Q is not used by this decision. The fixed-seven comparator
requests $\min(7,K)$ and applies the same geometric top-up, where $K$ is
the available common-worker pool size.
The separate online rule triggers at $k\geq5$ when three consecutive
complete outputs have equal node counts; otherwise it uses the full pool.
Equal counts do not establish equal geometry.

The quality panel contains five same-room groups, 21 images and five
buildings (5,376 replay attempts). Outcomes are first summarized
within each target view over its replay/history choices, then equally
over target views within group, groups within building, and buildings.
The compatible complete-time panel contains only three groups,
11 images and three buildings. Let $L_{vr}(k)$ be the recorded member
labor through prefix $k$, and $k_{7,vr}$ the fixed-seven-plus-geometry choice.
For a source/target split $(H,T)$, the project acquisition ratio is
\[
 \rho_r=
 \frac{\sum_{h\in H}L_{hr}(K)+
       \sum_{t\in T}L_{tr}(k_{{\rm stop},tr})}
      {\sum_{h\in H}L_{hr}(k_{7,hr})+
       \sum_{t\in T}L_{tr}(k_{7,tr})}.
\]
The archive averages these ratios over replays, splits, groups and buildings,
yielding 2.11526. This is an arithmetic mean of project ratios,
not a ratio of the target-only mean labor values
(1,660.3 versus 1,307.9 person-seconds), and not the geometric ratio in C8.
Replays reuse collected annotations and are not charged as repeated
real acquisitions. Unmeasured reference-production and assisted-connection
costs are not imputed as zero; recorded person-time is not wall-clock delay.

\paragraph{Source and implementation record.}
The source author is Sparkling-Flames (2026), \emph{3D\_Manhattan\_label},
a research software and output repository. All MR source IDs above resolve
to the full file titles, paths, Git blob identifiers and fixed-commit URLs
in the accompanying \texttt{source\_manifest.json} and
\texttt{SOURCE\_MAP.txt}, at commit
\texttt{a116cebbd03321f55c757130a4dd88b299539155}.
MR01--MR03 identify the profile definitions, result bindings and
underlying profile code; MR04--MR07 identify prediction scripts and matched
output tables; MR08--MR10 identify selection/stopping implementations,
common-support tables and cost summaries.
The source map links each method specification to its implementation and
reported output tables.


\subsection*{C10. Group-size precision and finite pools}
The 171-image main panel contains pools of 30--35 eligible annotators.
Table~\ref{tab:pool-capacity} separates participation fraction from absolute
group size. Each row uses a fixed-capacity image subset and its outcomes;
between-row differences also reflect the different images in each subset.

\begin{table}[ht]
\centering\small
\caption{Pool capacity, participation and 20-to-30-person changes on the fixed main panel. Q, P10 and standard deviation (SD) use scorable paths within each image. Complete-and-evaluable rates (CE) use all 256 attempts per image and size; their changes are in percentage points. Each column averages images equally.}
\label{tab:pool-capacity}
\setlength{\tabcolsep}{4pt}
\begin{tabular}{rrrrrrrr}
\toprule
Pool $N$ & Images & $20/N$ (\%) & $30/N$ (\%) & $\Delta Q$ & $\Delta$P10 & $\Delta$SD & $\Delta$CE (pp)\\
\midrule
30 & 12 & 66.7 & 100.0 & $-$0.911 & +1.612 & $-$3.003 & +1.107\\
31 & 9  & 64.5 & 96.8  & $-$1.620 & +2.834 & $-$4.966 & +0.087\\
32 & 8  & 62.5 & 93.8  & $-$0.973 & +1.482 & $-$1.080 & +1.221\\
33 & 22 & 60.6 & 90.9  & +1.697 & +4.598 & $-$3.430 & +0.621\\
34 & 22 & 58.8 & 88.2  & +1.541 & +4.293 & $-$1.814 & +0.444\\
35 & 98 & 57.1 & 85.7  & +0.486 & +2.396 & $-$1.454 & +0.251\\
\bottomrule
\end{tabular}
\end{table}

The precision checks support these aggregate changes. The approximate 95\% path-bootstrap interval for the 3--6-person mean gain was [4.629, 4.851]. For 20--30 people, intervals for mean Q, P10 and the within-image standard deviation were [0.439, 0.563], [2.620, 3.136] and [$-2.087$, $-1.962$], respectively. Building-level bootstrap intervals for the same late-stage changes were [0.037, 1.042], [1.942, 3.970] and [$-2.638$, $-1.503$]. Thus, later gains in lower-tail quality and reduced between-team variation persisted under both resampling schemes.

Among the 12 images with exactly 30 eligible annotators, mean Q fell by 0.911 points from 20 to 30 people, while P10 rose by 1.612; image-mean Q increased on nine images and decreased on three. Excluding these 12 images left 159 images. Their mean P10 gain from 20 to 30 people was 2.942 points, with increases on 155 images. Across all 171 images, footprint variation fell from 0.03592 to 0.01538; the estimated upper bound on the Monte Carlo standard error of this aggregate change was 0.000253. This bound combines paired estimates for 157 pools with sums of endpoint standard errors for the remaining 14. The precision assessment supports the main aggregate descriptions for the observed pools; individual-image low-tail estimates remain sensitive to discrete or multimodal team-quality distributions.



\subsection*{C11. Participant feedback and geometric outcomes}
The ordinary primary-target feedback subset contains 1,016 Manual responses on 118 images and 456 Semi responses on 34 images, with 24 people per condition. Occlusion was selected in 535 Manual and 175 Semi responses. These are multi-select reports from different task sets.

For Manual, the same-image omission difference between occlusion reporters and non-reporters was $-0.0481$ (95\% building-resampling interval [$-0.0747$, $-0.0124$]) across 72 images. The image-centered within-person difference was $-0.0418$ (95\% person-resampling interval [$-0.0624$, $-0.0208$]) across 24 people. The same-image simplification difference was $-3.75$ percentage points [$-6.41$, $-1.00$]; the within-person difference was $-2.08$ points [$-6.93$, $2.78$]. Intervals use 4,000 cluster resamples. Reports identify observed task conditions and can co-occur.

For Semi, the same-image omission difference was $+0.71$ percentage points (95\% building-resampling interval [$-0.27$, $+1.89$]) across 24 images in 11 buildings. The image-centered within-person difference was $-2.69$ points (95\% person-resampling interval [$-8.11$, $+0.76$]) across 24 people. These estimates show no consistent direction of association. Manual and Semi comparisons use their respective task sets.

On six expert-confirmed Semi tasks, the 138 retained ordinary responses comprise 20 unchanged submissions, one edited submission retaining a documented duplicate, 32 exact removals, and 85 edited submissions that did not cleanly resolve the target defect. These are a subset of the 156 original confirmed opportunities.

\subsection*{C12. Reconstructed annotation interface}
The interface links paired panorama endpoints with a three-dimensional preview. The reconstruction below also shows scope and difficulty-factor feedback, using archived field selections and labels from the saved interface configuration. Difficulty factors are participant-selected observations, distinct from the researcher-assigned easy, moderate, and difficult image groups.
\begin{figure}[ht]
\centering
\includegraphics[width=0.86\linewidth]{supplement_figures/interface_feedback.png}
\caption{Reconstructed annotation interface with scope and difficulty-factor feedback. The panorama, preview and displayed selections use the same archived annotation. This response records normal scope, occlusion, and stitching or stretching.}
\end{figure}
\FloatBarrier

\subsection*{C13. Quality and labor at two annotation budgets}
On 124 common contexts from 46 images and nine buildings, quality-selected teams had mean Q 93.835 at five people and 89.845 at thirteen; quality-and-time-selected teams had 94.055 and 89.747. The respective five-to-thirteen changes were $-$3.990 and $-$4.308 Q points. Recorded labor averaged 1,197 and 2,886 person-seconds for quality selection, versus 1,046 and 2,170 for quality-and-time selection. The geometric mean within-context labor growth ratios were 2.331 and 2.170, respectively; these are not ratios of the displayed arithmetic means. Teams were selected independently at each budget; the comparison therefore describes selection at two budgets.

\begin{table}[t]
\centering\small
\caption{Mean quality and within-context labor growth for independently selected five- and thirteen-person teams on 124 common contexts.}
\label{tab:budget-quality-labor}
\setlength{\tabcolsep}{3pt}\renewcommand{\arraystretch}{1.12}
\begin{tabular*}{\linewidth}{@{\extracolsep{\fill}}lrrr@{}}
\hline
Selection & Q, 5 & Q, 13 & Labor ratio\\\hline
Quality only & 93.835 & 89.845 & 2.331\\
Quality and time & 94.055 & 89.747 & 2.170\\\hline
\end{tabular*}
\end{table}


\subsection*{C14. Automatic and assisted outputs}

\paragraph{Same-room replacement directions.}
In paired-corner consensus, pairwise matching and MV50 support first fix the retained corner identities and their median coordinates. When team members retain different subsets of wall turns, no single member's submitted ring may cover every retained identity; fixed-node assisted ordering then supplies the cyclic connection order over those fixed nodes without adding or removing nodes, changing distinct-person support counts, or modifying median coordinates. The ten-person, five-person-replacement comparison has 322 specified events, 19 physical view pairs and seven buildings. Automatic source-ring connection makes 295 events jointly scorable; fixed-node assisted ordering expands this support to 318. Both branches yield positive pair-mean \(Q_C-Q_B\) in all 19 pairs and positive changes at both views in 18 pairs. Table~\ref{tab:automatic-same-room} compares the saved pair-level results under the same event--member-block--view-pair aggregation. On the original 295-event support, automatic and assisted effect summaries coincide to numerical precision.

\begin{table}[ht]
\centering\small
\caption{Same-room replacement effects from automatic and fixed-coordinate assisted outputs. The exchanged profiles, team sizes and 19 view pairs are shared; assisted ordering expands the jointly scorable event support.}
\label{tab:automatic-same-room}
\begin{tabular}{lrr}
\toprule
Outcome & Automatic & Assisted\\
\midrule
Jointly scorable events & 295/322 & 318/322\\
Mean \(G_S\) & 5.158 & 4.787\\
Mean \(Q_C-Q_B\) & 10.753 & 10.331\\
Positive pair-mean \(Q_C-Q_B\) & 19/19 & 19/19\\
Pairs positive at both views & 18/19 & 18/19\\
\bottomrule
\end{tabular}
\end{table}

The accompanying validation files retain the pair-level effects, support table and input specifications. The reported assisted means agree with the main comparison; the automatic branch establishes the same directional pattern on its own scorable support.

\paragraph{Assisted-order visibility diagnostics.}
The targeted 57-case connection-review queue contains 36 scorable outputs
with the camera in the declared footprint's visibility kernel, 20 scorable
outputs outside that kernel, and one unscorable self-intersection. The 20
flagged footprints are valid simple polygons containing the camera. The
kernel check tests the camera against all inward edge half-planes; it is a
geometric property of the declared polygon, not an image-based adjudication
of hidden structure. Conditional Q remains calculable for those 20 outputs.

Exact endpoint matching uses team membership and the assigned reference
identifier. None of these 57 endpoints enters the primary-target natural
paths for the 171-image main group-size panel.

Three flagged endpoints enter three scorable events in the 14-person,
seven-person-replacement ABC comparison. Eight flagged endpoints enter ten
events in the same-room ten-person, five-person-replacement comparison.
Table~\ref{tab:visibility-sensitivity} removes only those affected events and
reaggregates the saved scores with the original weights. Missing outcomes
are not assigned zero.

\begin{table}[ht]
\centering\small
\caption{Sensitivity to the flagged endpoints in the targeted assisted-order queue. The exclusion branch retains the same 69 images and 20 buildings for ABC, and 19 view pairs and seven buildings for the same-room comparison.}
\label{tab:visibility-sensitivity}
\setlength{\tabcolsep}{6pt}
\begin{tabular}{lrr}
\toprule
Outcome & Reported support & After exclusion\\
\midrule
ABC scorable events & 146 & 143\\
ABC mean $G_S$ & 3.441 & 3.686\\
ABC mean $Q_C-Q_B$ & 6.712 & 6.768\\
ABC images with $G_S>0$ & 50/69 & 49/69\\
ABC images with $Q_C-Q_B>0$ & 65/69 & 65/69\\
Same-room jointly scorable events & 318/322 & 308/322\\
Same-room mean $G_S$ & 4.787 & 5.072\\
Same-room mean $Q_C-Q_B$ & 10.331 & 10.527\\
Positive pair-mean $Q_C-Q_B$ & 19/19 & 19/19\\
Pairs positive in both views & 18/19 & 18/19\\
\bottomrule
\end{tabular}
\end{table}

ABC weighting proceeds through events, retained backgrounds, replacement
blocks, pools, and images. Same-room weighting proceeds through events,
member blocks, and physical view pairs, with the two views equally weighted
within a pair. The specified design counts are unchanged: 147 ABC
events and 322 same-room events. This sensitivity concerns the flagged endpoints in the 57-case queue.
Assistance duration was unrecorded for all 57 cases; recorded labor covers
the original annotations.

\clearpage
\subsection*{C15. Reconstruction, geometric states and branch criteria}

\paragraph{From panorama endpoints to geometry.}
The panorama has width $W$ and height $H_{\mathrm{img}}$, with pixel coordinates
measured from the upper left. Each upper--lower pair first receives the
midpoint of its shorter horizontal arc, denoted $u$; its original vertical
coordinates $v_t$ and $v_b$ are retained. In radians, let
$\theta=2\pi(u/W-1/2)$ and
$\beta_j=\pi(v_j/H_{\mathrm{img}}-1/2)$ for $j\in\{t,b\}$.
The camera is the origin, the vertical axis points upward, and camera height
is the length unit. For $0<\beta_b<\pi/2$, the horizontal radius is
$\rho=1/\tan\beta_b$. The floor point is
$(\rho\sin\theta,-1,-\rho\cos\theta)$ and the corresponding upper point is
$(\rho\sin\theta,-\rho\tan\beta_t,-\rho\cos\theta)$.
Thus the corner's wall height is $h=1-\rho\tan\beta_t$.
The stored endpoint representation uses $W=1024$ and $H_{\mathrm{img}}=512$.
Measurements in metres, such as the depth example in the main paper, apply
the stated camera-height multiplier to these normalized coordinates.

\paragraph{Geometric validity and output states.}
The scoring geometry is an implicitly closed paired wall ring. Checks require
finite coordinates, equal upper and lower ring sizes, at least three pairs,
vertical pairing, the common floor plane, strictly positive wall heights,
and a simple nondegenerate footprint. Zero-length edges, self-intersections
and unresolved numerical degeneracies retain their failure states.
The spherical-boundary measurements additionally require defined finite
projections and nondegenerate shorter arcs. A valid floor can therefore
exist without a complete scorable wall layout. The scorer measures internal
orthogonal regularity and local height variation on this geometry; it does
not fit them away. Camera position and visibility-kernel diagnostics are
recorded separately, as in C14.

Scale-dependent numerical guards handle degenerate geometry and undefined projections; their implementation is specified in the accompanying source map.

\begingroup\small
\begin{longtable}{@{}>{\raggedright\arraybackslash}p{.24\linewidth}>{\raggedright\arraybackslash}p{.14\linewidth}>{\raggedright\arraybackslash}p{.56\linewidth}@{}}
\toprule
Branch & Minimum nodes & Additional condition and output\\
\midrule\endfirsthead
\toprule Branch & Minimum nodes & Additional condition and output\\\midrule\endhead
Geometric Q & 3 paired corners & Both assessed and reference layouts pass the geometric and metric-domain checks above. Valid zero floor overlap gives Q=0; unavailable geometry retains an unscorable status.\\[3pt]
Automatic consensus order & 3 selected identities & A projected source ring covers all selected identities with defined positions. Its connected geometry then undergoes the Q checks; a candidate order alone does not establish scorability.\\[3pt]
Fixed-coordinate assisted order & Same selected set & The supplied order preserves identities, support and coordinates. The resulting ring undergoes the same validity checks.\\[3pt]
Natural-size complete-output rates & 4 selected identities & A valid footprint gives a complete floor; paired wall geometry with recorded status \texttt{ok} gives the complete-evaluable geometry flag. Q additionally requires successful metric calculation.\\[3pt]
Source-history stopping & 4 nodes in a complete output & Starting at the requested prefix, add members until the recorded complete-geometry criterion is satisfied. Target Q does not choose the stopping point.\\
\bottomrule
\end{longtable}
\endgroup

These minimums belong to their respective procedures. The general geometry
and source-ring construction admit three nodes; the natural-size completion
flags and stopping delivery rule impose four.

\paragraph{Reference self-occlusion.}
For a valid reference footprint containing the camera strictly inside,
$H$ (\texttt{ref\_H}) counts floor vertices whose open sightline from the
camera crosses the interior of a non-incident footprint edge before reaching
the vertex. Tangential vertex contacts and collinear contacts are recorded
separately from crossings. Invalid footprints or a camera outside the
interior leave this feature unavailable. This is concealment by the room
boundary in the floor model. The same count enters the opportunity features
in C6 and the matched cross-view comparisons.

The accompanying source map in
\texttt{supplement\_sources/implementation\_reference/} identifies the
reconstruction, validity, scoring, correspondence and visibility functions
used to check these definitions.

% Include after sections/supplement_c_geometry in supplement_personnel.tex.
\clearpage
\subsection*{C16. Parameter, validation and implementation details}
\label{sec:consensus-implementation}

\begin{figure}[ht]
\centering
\includegraphics[width=\textwidth]{supplement_figures/geometry_consensus.pdf}
\caption{From paired-corner observations to collective geometry and quality. Distinct-person support determines retained identities and median positions. Automatic or assisted ordering connects these nodes, and valid layouts are assessed against the matched target reference. The outputs retain structural support, formation status and geometric quality.}
\label{fig:geometry-consensus}
\end{figure}

\paragraph{Quality-score development.}
Parameter development used controlled geometry and development cases. Twenty-eight configurations compared direction scales of approximately 6.12, 8, 10 and 12 degrees, boundary denominators of 32, 48 and 64, and combination exponents of 1.25, 1.5, 2 and 3. The retained \(8^\circ\) scale responds to isolated direction deviations; L2 moderates their accumulation with footprint loss, and the boundary denominator of 32 retains sensitivity to omissions and localization errors. Development comparisons of global-height scales 1.0 and 0.8 and local-height settings near 0.25 informed the retained height scales. The formula was then fixed for validation and personnel analyses.

\paragraph{Score-scale sensitivity.}
The sensitivity analysis uses nine distinct parameter settings: the primary score and eight variants multiplying one scale by $0.75$ or $1.25$. The $1.25\times$ direction variant is the $10^\circ$ setting. Boundary variants scale $\sqrt{32}$, giving squared denominators of 18 and 50; global-height scales are 0.6 and 1.0, local-height scales 0.1875 and 0.3125, and direction scales $6^\circ$ and $10^\circ$.

The 3--6-person comparison uses 206 images and the 20--30-person comparison uses 171. Each interval retains the same jointly scorable paths across settings within each image; P10 changes subtract endpoint quantiles. The 14-person A--B--C comparison uses 69 images in 20 buildings, includes fixed-node assisted ordering and retains the aggregation defined in C17. Every setting preserves early mean gains, late P10 gains, reduced late standard deviation and positive replacement excess (Table~\ref{tab:score-scale-sensitivity}). The existing parameter and result records accompany the manuscript in \path{supplement_sources/score_sensitivity/}.

\begin{table}[htbp]
\centering\small
\caption{Score-scale sensitivity on fixed output supports. All values are Q-point changes or Q-point excesses. Each row is a distinct parameter setting.}
\label{tab:score-scale-sensitivity}
\setlength{\tabcolsep}{4pt}
\begin{tabular}{@{}lrrrrr@{}}
\toprule
Setting & Mean 3--6 & Mean 20--30 & P10 20--30 & SD 20--30 & $G_S$ \\
\midrule
Primary & +4.738 & +0.502 & +2.844 & -2.005 & +3.441 \\
Global height $\times0.75$ & +4.743 & +0.503 & +2.844 & -2.003 & +3.463 \\
Global height $\times1.25$ & +4.735 & +0.502 & +2.846 & -2.006 & +3.430 \\
Local height $\times0.75$ & +4.832 & +0.515 & +2.865 & -2.009 & +3.455 \\
Local height $\times1.25$ & +4.692 & +0.497 & +2.835 & -2.003 & +3.432 \\
Direction $\times0.75$ & +5.223 & +0.679 & +3.251 & -2.205 & +3.460 \\
Direction $\times1.25$ ($10^\circ$) & +4.374 & +0.400 & +2.597 & -1.880 & +3.385 \\
Boundary scale $\times0.75$ & +5.011 & +0.510 & +2.999 & -2.104 & +3.989 \\
Boundary scale $\times1.25$ & +4.600 & +0.505 & +2.784 & -1.955 & +3.158 \\
\bottomrule
\end{tabular}
\end{table}


\paragraph{Eight-reviewer questionnaire protocol and defect relations.}
Two Manual annotations were randomly selected from each of 30 panoramas across 17 buildings: nine easy, nine moderate and nine difficult ordinary panoramas plus three doorway panoramas. Eight annotators were randomly invited to take part in the questionnaire. They rated each annotation as good, medium or poor and selected visible defects, using only panorama overlays, shared-scale floor footprints and three-dimensional views in a fixed randomized arrangement. Grades were encoded as 2, 1 and 0, respectively, before averaging.

Of 480 submitted grades, 478 were valid (67 good, 163 medium and 248 poor); two were unknown. The 238 complete reviewer--image pairs comprise 123 ties and 115 directional comparisons. Q agrees with 110 directions and opposes five; mean-height prism IoU agrees with 88 and opposes 27. Neither metric ties on these 115 directional pairs. Their correlations with mean grades across the 60 annotations are 0.946 and 0.738. Averaging the common reviewers within each image gives 21 non-tied comparisons, with agreement in 20 for Q and 16 for IoU. These grade-based comparisons complement the defect feedback.

Defect-selection shares use valid answers for each annotation: unknown and blank answers are excluded, while ``none'' is a valid report of no visible defect. The 480 defect answers comprise 460 valid reports and 20 unknown reports, with no blank answers or combinations of ``none'' with a specific defect. Across the 60 annotations, missing-structure shares correlate with missing effective turns (\(\rho=0.544\)) and floor omission (\(\rho=0.240\)); explicit floor-omission shares correlate with omitted floor fraction (\(\rho=0.586\)). Extra-structure shares correlate with extra turns (\(\rho=0.415\)) and floor extension (\(\rho=0.157\)). Boundary-position shares correlate with \(S_b\) (\(\rho=0.395\)) and \(D_{\mathrm{dir}}\) (\(\rho=0.358\)). Wall-height shares correlate with \(H_m\) (\(\rho=0.090\)), \(H_c\) (\(\rho=0.612\)) and the combined height loss (\(\rho=0.609\)), distinguishing local height discrepancies from an overall height shift.

Reviewers assigned the same grade in 1,369 of 1,666 valid within-annotation reviewer pairs (82.2\%). Resampling the 17 buildings 10,000 times, retaining all their images, paired annotations and the eight reviewers, gives percentile 95\% intervals of [0.870, 0.977] for Q's correlation and [0.496, 0.880] for prism IoU. Their correlation difference is 0.208, with interval [0.062, 0.440]. Omitting each reviewer in turn gives Q correlations of 0.900--0.960 and prism-IoU correlations of 0.695--0.760.

The questionnaire responses, per-annotation geometric measurements and recomputation scripts are provided in \path{supplement_sources/questionnaire/}.

\paragraph{Additional geometric review.}
A separate geometric review examined 25 same-image pairs: 23 Manual and two Semi. Twenty pairs were selected for specified geometric conditions and score relationships; five mechanism examples were already familiar to the reviewer. Across all 25 pairs, ten received directional judgments, three were judged close, and 12 were difficult to compare or problematic on both sides. Among the 20 condition-selected pairs, six judgments were directional. Q with direction scales \(8^\circ\) and \(10^\circ\) agreed with all five directional Manual judgments, while prism IoU agreed with three. In the directional Semi pair, the reviewer preferred the lower-scoring response.

\paragraph{Working skip penalty.}
The matching uses a fixed penalty of \(g=7.5^\circ\) for each unmatched observation. Selected during development, this value balances positional variation among observations of the same turn against the separation of nearby distinct turns. Its effect operates through the cyclic partial matching defined in the main text.

\paragraph{Candidate identity groups.}
Pairwise cyclic alignment uses the average spherical distance of the upper
and lower endpoints and a skip penalty of \(g=7.5^\circ\) for each unmatched
observation. Complete-linkage agglomerative clustering encodes these
alignments with dissimilarity zero on the diagonal, \(c(a,b)/180^\circ\) for
observations selected as a pair in the corresponding worker-pair alignment,
and 2 otherwise. The hierarchy is cut at 1. Each resulting candidate group
therefore contains at most one observation per worker, with matching support
for every pair of observations. Pairwise links through an intermediate
observation alone are insufficient to join a group.

\paragraph{Periodic positioning and ties.}
The observed-pair anchor minimizes the sum of the larger upper- or
lower-endpoint angular distances to the other group members. Anchor ties
follow the fixed worker, response and pair-index order. With panorama
width \(W=1024\), anchor coordinate \(x_0\), and wrapped offsets
\(d_j=((x_j-x_0+W/2)\bmod W)-W/2\), a position is returned only if
\(\max_j d_j-\min_j d_j<W/2\) and no \(|d_j|=W/2\). The boundary checks
use a \(10^{-9}\)-pixel tolerance. A failed check leaves the position
unresolved. Horizontal and vertical medians use the same group's
observations with equal worker weights; their resulting coordinates may
lie between original observations.

\paragraph{Connection of retained corners.}

Corner ordering enters at two distinct stages: individual annotation and collective connection. During Manual annotation, each participant places corner pairs and orders their own closed wall ring in the viewer to specify the sequence of adjacent walls. After pairwise matching and MV50 selection, the team's retained identities \(\mathcal V\) and their median coordinates are already determined, but individual members may have retained different subsets of wall turns.

Automatic connection tests whether the team's submitted rings already specify a complete order over \(\mathcal V\). Let \(\rho_w|_{\mathcal V}\) be worker \(w\)'s cyclic order restricted to \(\mathcal V\). For \(|\mathcal V|\geq3\), \(\mathcal R(\mathcal V)\) contains those projected rings that cover every selected identity in \(\mathcal V\) exactly once, treating cyclic shifts and whole-ring reversal as equivalent; with fewer selected identities, the set is empty. When \(\mathcal R(\mathcal V)\) is nonempty, the automatic branch selects the most-supported projected source ring:
\begin{equation}
 \rho^*=\underset{\rho\in\mathcal R(\mathcal V)}{\arg\max}
 \sum_{w\in S}\mathrm{1}\!\left[
 \operatorname{canon}(\rho_w|_{\mathcal V})=\rho\right].
 \label{eq:ring}
\end{equation}
When no single submitted ring covers all retained identities in \(\mathcal V\)---or when a reviewed connection order is specified---the retained corner identities still have defined geometric positions in the panorama, and their cyclic connection order is supplied by fixed-node assisted ordering. This step performs the same wall-sequence ordering operation used during Manual annotation, now applied to the team's fixed consensus nodes: it connects the retained identities in \(\mathcal V\) without adding or removing nodes, altering distinct-person support counts, or modifying median coordinates. Both automatic and assisted outputs then undergo identical 3D reconstruction and geometric validity checks, retaining explicit branch labels and failure states (fewer than three nodes, unresolved order, or self-intersecting footprint; Supplement C15).


\paragraph{Source rings and output records.}
Cyclic shifts and whole-ring reversal are canonicalized as equivalent.
Ties between the most-supported covering source rings follow a fixed
canonical order and are recorded. A direct adjacency occurs in an original
source ring; a projected adjacency results from removing intervening
unselected nodes. Records retain the supporting workers for both kinds of
adjacency. Candidate rings, fixed-coordinate assisted orders and geometric
validity states are stored separately.

Each team record retains membership, original responses, pairwise matches,
identity groups and source pairs, including identities below the MV50
threshold. Checks recount distinct supporters, verify the at-least-half
selection, recompute periodic and vertical medians, and verify selected-node
coverage and ring support. Thus each retained corner can be traced to its
source observations, supporter count, median position and place in the final ring.

\paragraph{Floor-region comparison and local structure.}
BEV voting partitions the team's overlapping floor polygons into regions and retains regions covered by at least half the people. Its output is a two-dimensional floor region. Paired-corner consensus retains the upper and lower coordinates of supported wall turns and connects them into a wall layout.

Figure~\ref{fig:consensus-unb45} shows the full panorama and the corresponding floor outputs for a ten-person team. The same seven people mark E and F. Both the paired-corner layout and the BEV region retain the resulting local change in floor coverage; the paired-corner output also records the supporters and upper--lower positions of those turns.

\begin{figure}[htbp]
\centering
\includegraphics[width=\linewidth]{supplement_figures/consensus_unb45.pdf}
\caption{Local wall turns enter a collective layout through their supporters. In panorama uNb9QFRL6hY-45, E and F each receive support from the same seven of ten annotators. The paired-corner layout retains six turns, while BEV voting retains the corresponding floor-region change. The panorama and close view locate E and F; floor views share coordinates and scale, with the reference shown by dashed lines.}
\label{fig:consensus-unb45}
\end{figure}

\FloatBarrier
\subsection*{C17. Recombination, target assignment and outcome estimands}

\paragraph{Target-compatible pools.}
Each response is assigned to its actual target before evaluation. Where available, manual target adjudication takes precedence. Automatic assignment considers distinguishing wall segments, turns and exclusive regions, then uses floor IoU and normalized boundary distance to resolve ties. Assignment uses the complete submitted geometry and excludes the resulting composite Q. A pool is defined by image, annotation condition and target/reference. References for admitted alternative spaces combine the adjudicated floor extent with the confirmed image height.

\paragraph{Profile and member comparisons.}
The main text introduces the structural-expression, consensus-response and joint profiles. Supplement C6 specifies the feature definitions and fitting procedure. The designs below distinguish fixed-size composition, whole-class combinations, replacement scale and retained backgrounds. The three-team design has equal-size, mutually disjoint replacement blocks and a common retained background. Direct replacement fixes the background, team size and replacement size. The bidirectional six-person comparison requires all four outputs scorable on the same images.

\begin{table}[ht]
\centering\small
\caption{Controlled team comparisons by fixed conditions and varied factors.}
\label{tab:comparison-designs}
\setlength{\tabcolsep}{5pt}\renewcommand{\arraystretch}{1.13}
\begin{tabular}{@{}p{.20\textwidth}p{.35\textwidth}p{.40\textwidth}@{}}
\hline
Comparison & Held fixed & Varied factor / question\\\hline
Matched replacement & Image, target, team size, retained members & Same- versus cross-profile blocks; magnitude and direction\\
Direct replacement & Image, target, team size, retained members & One exchanged block; directional change\\
Fixed-size composition & Image, target, background, candidate pool & Profile counts; members at identical counts\\
Whole-class combination & Image, target, selected profile definitions & All members of each selected class; team size varies\\
Replacement scale & Starting team and target & Number of replaced members\\
Retained background & Image, target, exchanged blocks & Background members and their number\\
Cross-view comparison & Physical-room relation and member blocks & Separate original annotations at each viewpoint\\\hline
\end{tabular}
\end{table}

\paragraph{Construction of the matched replacement events.}
The 14-person, seven-person comparison uses 147 specified events on 69 images in 20 buildings. Replacement blocks were constructed within eligible image--target pools, with mutually disjoint seven-person blocks X and Y from C2, Z from C3, and a seven-person retained block D. Registration reused existing blocks after checking their current class eligibility and disjointness, and supplemented them with seeded member draws (base seed 2026102907). Backgrounds were drawn from the remaining people at feasible low, middle and high counts of the incoming class, or at one feasible stratum selected by a deterministic design key. The focused registration allowed two replacement blocks per design cell; the screening registration allowed one, with one background per selected stratum. Duplicate events were removed. This specified event set supplies the matched comparisons; 146 events have all three outputs scorable.

\paragraph{Natural and prioritized entry orders.}
The principal natural-order analysis contains 223 Manual pools, with 256 random entry orders per pool. Each retains its assigned target, reference and eligible annotators. The main fixed panel uses 171 ordinary images with at least 30 people. One person supplies an individual response, two-person aggregation is calculated separately, and natural group curves principally use \(k\geq3\). Fully enumerated finite team sets give exact marginal distributions; pathwise differences retain observations along the same order. Historical-priority comparisons use paired base orders, principally 32 pooled paths, with priorities based on historical Q, richer turning structure, shorter recorded time or the largest structural class.

\paragraph{Natural-size estimands and uncertainty.}
To make the natural-size estimands explicit, fix an image and target pool with \(N\) eligible people. For a uniformly random entry order \(\pi_r\), let \(S_{rk}=\{\pi_r(1),\ldots,\pi_r(k)\}\). For any particular \(k\)-person subset \(A\), \(\Pr(S_{rk}=A)=\binom Nk^{-1}\). Let \(M\) be the number of sampled orders, \(\mathcal J_k\) the indices with a scorable output at size \(k\), and \(q_{rk}\) their quality scores. Let \(e_{rk}=1\) when output \((r,k)\) is complete and evaluable, and zero otherwise. The conditional mean quality and complete-and-evaluable rate are
\begin{equation}
 \begin{aligned}
 \widehat\mu_Q(k)&=\frac{1}{|\mathcal J_k|}\sum_{r\in\mathcal J_k}q_{rk},\\
 \widehat p_{\mathrm{ce}}(k)&=\frac1M\sum_{r=1}^{M}e_{rk}.
 \end{aligned}
 \label{eq:natural-outcomes}
\end{equation}
Order indices retain duplicate-team occurrences. If \(\mathcal J_k\) is empty, mean quality is undefined. For two sizes, the paired pathwise change instead uses \(\mathcal J_{k,l}=\mathcal J_k\cap\mathcal J_l\):
\begin{equation}
 \widehat\Delta_{k,l}^{\mathrm{paired}}
 =\frac{1}{|\mathcal J_{k,l}|}\sum_{r\in\mathcal J_{k,l}}
 (q_{rl}-q_{rk}).
 \label{eq:paired-size-change}
\end{equation}
It is undefined without jointly scorable paths and need not equal \(\widehat\mu_Q(l)-\widehat\mu_Q(k)\). These quantities are first calculated within each image's analysis pool; reported fixed-panel curves average the image-level statistics over the stated common image set.


Monte Carlo precision was assessed by 2,000 bootstrap resamples of complete entry paths within each image pool, preserving the original generation-batch sizes. Each resample used the same path weights across group sizes. Q summaries retained their scorable denominators, whereas completion used all attempts. Image-specific P10 values were computed by linear interpolation before equal-weight averaging across images. Approximate 95\% percentile intervals describe path-sampling uncertainty conditional on the observed pools. Separate building-level bootstrap intervals resampled buildings while retaining the path point estimates, assessing variation across the 22 sampled buildings.

\paragraph{Variance across compositions and member sets.}
For panel \(p\), let \(\mathcal T_{pc}\) be the scorable teams with profile-count vector \(c\), \(n_{pc}=|\mathcal T_{pc}|\), and \(n_p=\sum_c n_{pc}\). Write \(\overline q_{pc}\) for the composition mean and \(\overline q_p\) for the mean over all scorable teams. The finite-population variances are
\begin{equation}
 \begin{aligned}
 V_{p,\mathrm{within}}
 &=\frac1{n_p}\sum_c\sum_{t\in\mathcal T_{pc}}
 (q_t-\overline q_{pc})^2,\\
 V_{p,\mathrm{between}}
 &=\sum_c\frac{n_{pc}}{n_p}(\overline q_{pc}-\overline q_p)^2,\\
 V_{p,\mathrm{total}}
 &=V_{p,\mathrm{within}}+V_{p,\mathrm{between}}.
 \end{aligned}
 \label{eq:composition-variance}
\end{equation}
Compositions receive weights proportional to their scorable team counts. Variance components are then averaged over panels within target pool, pools within image, images within building, and buildings equally. With normalized panel weight \(a_p\), the between-composition share is the ratio of the aggregated variances:
\begin{equation}
 \eta_{\mathrm{between}}
 =\frac{\sum_p a_pV_{p,\mathrm{between}}}
 {\sum_p a_pV_{p,\mathrm{total}}}.
 \label{eq:composition-share}
\end{equation}
Panels with no scorable teams have undefined variances. Zero-variance panels retain zero components, and a zero aggregate total leaves the share undefined. The reported share describes the scorable finite team sets under these weights.


\FloatBarrier
\subsection*{C18. Complete measurement and positioning definitions}

\paragraph{Continuous floor-boundary distance.}
Let \(\mathcal{B}_P,\mathcal{B}_R\) contain midpoint samples at equal-length intervals on the spherical floor boundaries of layouts \(P,R\). Let \(\Gamma_P,\Gamma_R\) be their continuous shorter great-circle arcs between consecutive corners. With spherical angular distance \(\theta\) in degrees, the symmetric distance is
\begin{equation}
 \begin{aligned}
 2S_b(P,R)={}&\frac{1}{|\mathcal B_P|}\sum_{p\in\mathcal B_P}
 \min_{r\in\Gamma_R}\theta(p,r)\\
 &+\frac{1}{|\mathcal B_R|}\sum_{r\in\mathcal B_R}
 \min_{p\in\Gamma_P}\theta(r,p).
 \end{aligned}
 \label{eq:supp-boundary-distance}
\end{equation}
Every sample is compared with the other layout's continuous boundary. The implementation uses 2,048 samples in each direction and gives the two directed arithmetic means equal weight.

\paragraph{Global and local height.}
Heights interpolate linearly along floor edges. For cyclic corner heights \(h_i\) and floor-edge lengths \(\ell_i\), the perimeter-weighted mean height and relative mean-height error are
\begin{equation}
 \mu_P=\frac{\sum_i\ell_i(h_i+h_{i+1})/2}{\sum_i\ell_i},
 \qquad H_m(P,R)=\frac{|\mu_P-\mu_R|}{\mu_R}.
 \label{eq:supp-height-mean}
\end{equation}
Subtract each layout's own mean to obtain \(\widetilde h_P(x)=h_P(x)-\mu_P\) on \(\partial A_P\), and similarly for \(R\). For midpoint samples \(\mathcal X_P\) at equal floor-perimeter intervals, define the nearest horizontal-floor positions
\begin{equation}
 \mathcal N_R(x)=\underset{y\in\partial A_R}{\arg\min}\,\|x-y\|_2.
 \label{eq:supp-nearest-floor-set}
\end{equation}
The directed centered-height error is
\begin{equation}
 \operatorname{MSE}^{c}_{P\rightarrow R}
 =\frac{1}{|\mathcal X_P|}\sum_{x\in\mathcal X_P}
 \max_{y\in\mathcal N_R(x)}
 [\widetilde h_P(x)-\widetilde h_R(y)]^2.
 \label{eq:supp-height-directed}
\end{equation}
The maximum resolves ties in nearest floor position. Exchanging \(P,R\) gives the reverse error, and
\begin{equation}
 H_c(P,R)=\frac1{\mu_R}
 \sqrt{\frac{\operatorname{MSE}^{c}_{P\rightarrow R}
 +\operatorname{MSE}^{c}_{R\rightarrow P}}2}.
 \label{eq:supp-height-center}
\end{equation}
The calculation uses 4,096 samples per direction. Both height errors are dimensionless; a changed corner can affect the global and local components.

\paragraph{Score scales and geometric comparisons.}
The boundary-loss denominator is \(32\,\mathrm{deg}^2\), the direction scale is \(8^\circ\), and the direction loss approaches 1.5. Global and local height scales are 0.8 and 0.25. Thus \(S_b=\sqrt{32}^\circ\) contributes unit boundary loss, an \(8^\circ\) direction residual contributes 0.75, and each height loss equals one at its stated scale. The L2 combination lies between the maximum and sum of the floor and direction losses. Q reaches 100 when the footprint and height profiles agree with the reference and the assessed layout has zero internal direction residual.

The controlled examples distinguish local height, structural completeness and geometric deviation. In the square-room control with reference height 2.6 camera heights, raising one side by 20\% and lowering the opposite side by 20\% preserves the footprint and mean height. Prism IoU remains 1, \(H_m=0\), \(H_c=0.1633\), and Q falls from 100 to 70.09. Column-wise layout-depth RMSE is 0.4437\,m under the 1.6\,m camera-height convention. In the 3\,m controls, raising one corner by 2.4\,m or every corner by 0.6\,m gives the same corner mean and prism IoU of 83.33\%; Q is 46.97 and 94.12, respectively. Omitting the deepest tested narrow recess removes four turns, with prism IoU at 97.09\% and Q at 96.67. Turn counts record this structural loss alongside geometric scores (Supplement A).

A real team-endpoint comparison also separates overlap, structure and boundary accuracy. Floor IoU changes from 0.8414 to 0.8451 and Q from 75.03 to 78.80. Bottom-corner recall at \(9^\circ\) rises from 0.5 to 1, while layout-depth RMSE falls from 0.2359\,m to 0.1851\,m. These component measurements describe simultaneous changes in the complete geometry.

\paragraph{Support rule.}
For candidate group \(G\) in team \(S\), distinct-person support is
\begin{equation}
 s(G)=|\{w\in S:\exists\,a_{wj}\in G\}|.
 \label{eq:supp-support}
\end{equation}
For fixed \(G\), the binary decision \(z\in\{0,1\}\) incurs disagreement
\begin{equation}
 \mathcal L_G(z)=s(G)(1-z)+[|S|-s(G)]z.
 \label{eq:supp-support-disagreement}
\end{equation}
The MV50 rule \(z=\mathrm{1}[2s(G)\geq|S|]\) minimizes this disagreement, retaining the corner when the two alternatives tie.

\paragraph{Periodic position.}
The unwrapping anchor is an observed pair minimizing the sum of the larger endpoint distances:
\begin{equation}
 a^*=\underset{a\in G}{\arg\min}\sum_{b\in G}
 \max\{\theta(p_a^t,p_b^t),\theta(p_a^b,p_b^b)\}.
 \label{eq:supp-position-anchor}
\end{equation}
Let \(x_0\) be its aligned horizontal coordinate, \(W=1024\) the panorama width, and \(x_j\) the aligned source coordinates. With \(\operatorname{wrap}_W(d)=((d+W/2)\bmod W)-W/2\), the retained corner coordinates are
\begin{align}
 x_G&=\left[x_0+\operatorname{median}_j\{\operatorname{wrap}_W(x_j-x_0)\}\right]\bmod W,\label{eq:supp-median-x}\\
 y_G^t&=\operatorname{median}_j y_j^t,\qquad y_G^b=\operatorname{median}_j y_j^b.\label{eq:supp-median-y}
\end{align}
All three medians use the same supporting people with equal weights. The spread must fit within an open half-panorama; C16 specifies this check, tie resolution and the source-ring connection procedure.

\FloatBarrier
\subsection*{C19. Extended individual, team and cross-view results}

\paragraph{Individual geometric and labor responses.}
On 201 images shared by C1, C2 and C3, mean individual Q was 64.965, 66.913 and 80.971, respectively. Mean matched-boundary errors were $1.856^\circ$, $1.730^\circ$ and $0.941^\circ$. This error averages absolute upper- and lower-boundary angular deviations at matched corner pairs and differs from footprint boundary distance $S_b$. On the 130 easy images shared by C2 and C3, their turn-count difference was 0.058 while their Q difference was 13.943 points.

Within-person, within-building difficult-minus-easy comparisons gave C3 an additional 2.330 turns relative to C2 across fifteen buildings, positive in all fifteen. The corresponding Q difference was +5.235 points, positive in nine buildings. Adding geometric features to the difficulty, annotator and source controls reduced held-out-building turn-deficit MSE from 1.744 to 1.163. Reference turn count is among those predictors and also enters the definition of turn deficit; binary simplification MSE and Q MAE are reported alongside it in the main text. Supplement C7 specifies the prediction procedure.

In 100 slower--faster pairs reusing 35 annotators, the slower person's easy-to-difficult advantage changed by $-$0.054 turns and $-$0.029 short edges on average. Positive changes occurred in 53 and 55 pairs, respectively. Across 162 images and seventeen buildings, J2 minus J1 gave Q $-$0.133 and a geometric mean labor ratio of 0.885. These outcomes describe recorded person-time and structural expression.

The six expert-confirmed controlled Semi tasks yielded 156 opportunities from 26 annotators, with 25 unchanged submissions. The retained ordinary primary-target subset has 138 responses and the issue-specific outcome counts in C11. Three separate natural-initialization submissions for e9zR4mvMWw7-39 also retained their initial coordinates. The natural initialization had been checked by experts before the experiment. Supplement B20 shows the natural example. In the duplicated-pair example in B8, panorama Z6MFQCViBuw-10, A2 removed the added pair while retaining every other coordinate and A3 retained it.

\paragraph{Small groups, completion and retrospective stability.}
On the fixed 171-image panel, the consecutive changes in mean Q for sizes 3--6, 6--10, 10--20 and 20--30 were +4.74, +1.75, +1.31 and +0.50 after rounding. The corresponding P10 gains were +8.27, +3.73, +3.78 and +2.85; each difference was calculated before rounding the displayed endpoint levels. Actual failure rates rose by 1.894 percentage points from three to six people and fell by 0.416 points from twenty to thirty. Complete-and-evaluable rates changed by $-$1.857 and +0.420 points. An output can be produced but unscorable, accounting for the distinction between these rates.

Among 98 images with capacity 35, 97 had scorable outputs at both thirty and 35 people. On these 97 images, mean Q increased by 0.281 points, P10 by 2.736 and within-image standard deviation fell by 1.952. Complete capacity strata and the full-pool diagnostics are given in C10.

Retrospective stability requires the range of image-mean Q from a candidate size through the actual full pool to be at most a specified tolerance, at least $\max(5,\lceil0.15N\rceil)$ additional members to remain observable, and every size to have at least 64 scorable paths. With a two-point tolerance, onsets were detected on 163 of 206 ordinary images: 123/130 easy, 25/39 moderate and 15/37 difficult images. The remaining 43 had no detected onset. Median sizes among detected onsets were nine, sixteen and seventeen. Tolerances of one, two, three and five points yielded 145, 163, 170 and 184 detected onsets, respectively. The criterion uses the complete observed curve.

Early-entrant composition tracked later quality within the stored natural paths. Orders beginning with six C2+C3 entrants averaged about 1.78 Q points higher over sizes 6--20 than those beginning with C1+C2 entrants, with positive differences on 165 of 190 comparable images. These orders vary in both profiles and particular people. Separately calculated one-to-two and two-to-three paired changes averaged about +4.23 and +3.61 Q points on the 171-image panel, requiring both endpoints scorable for each contrast. At two people, either person's support meets the at-least-half rule.

\paragraph{Additional composition and replacement comparisons.}
Increasing the C3 count by one at the expense of C2 in fixed 20-person groups raised composition-mean Q by 1.192 points, positive on 27 of 35 images. This comparison averages member sets within a composition. On common quality/time support for twelve-person teams, C2-to-C3 gave +1.963 Q and a geometric labor ratio of 1.020 on 27 images; J3-to-J4 gave +2.701 Q and 0.897 on fourteen images.

The direct automatic 20-person, seven-person C2-to-C3 replacement increased Q by 4.484 points, positive on 77 of 84 images. On its forty images with joint quality and time support, the increase was 4.406 Q points and 121.05 person-seconds. For the 14-person A--B--C comparison, the automatic-only support comprises 67 images in twenty buildings: $G_S=3.767$, positive on 49 images. The main comparison including assisted ordering has 69 images, $G_S=3.441$ and fifty positive image effects.

\begin{table}[ht]
\centering\small
\caption{Bidirectional six-person comparisons replacing three people. Both directions require all four outputs scorable on the same 36 images in ten buildings.}
\label{tab:bidirectional-half}
\begin{tabular}{lrrr}
\toprule
Direction & Symmetric excess $G_S$ & Positive excess & $Q_C-Q_B$\\
\midrule
C1 to C3 & 3.269 & 32/36 & +7.327\\
C3 to C1 & 3.955 & 24/36 & $-$7.777\\
\bottomrule
\end{tabular}
\end{table}
On this shared support, the C3-versus-C1 team-type mean difference is +7.865 Q points. It compares team types across the same images, separately from the two directional exchanges in the table.

The R1/R2 consensus-response profiles also differed in individual geometry. Across 192 common images, R2 exceeded R1 by 13.937 Q points, positive on 175; its matched-boundary error was $0.849^\circ$ lower and its omitted-area fraction 0.0545 lower. In the twelve-person, four-person matched replacement, within-profile Q-change magnitude averaged 4.132 and cross-profile magnitude 6.308, with the latter larger on 62 of 83 images.

\paragraph{Retained backgrounds and geometry.}
In the joint-profile comparison holding seven-person blocks fixed, increasing total size from twenty to 28 reduced absolute Q change by 1.317 points on 76 images, with reductions on 63. All four outputs were scorable on each path; the mean difference in signed effects between sizes was +0.014 points.

The background-reversal comparison fixes image, reference, team size, replacement size and exchanged members, using unchanged automatic endpoints. Of 276 blocks observed under at least two backgrounds, 56 reversed sign across seven images; 24 crossed both +1 and $-$1 Q point across five images. Requiring identical background profile counts yielded 120 subgroups, representing 108 distinct blocks across seven images and four buildings. Eighteen subgroups reversed sign; nine crossed both thresholds across four images. Blocks have differing numbers of observed backgrounds. The q9vSo1VnCiC-30 example in the main text had the largest effect range among comparisons with the same exchanged members and background profile counts.

Score range and geometric components explain additional response differences. In the automatic 20-person, seven-person C2-to-C3 comparison, 1,157 pairs with both Q values at least ninety differed by 2.064 points on average; 672 pairs with at least one lower value differed by 11.012. An image can contribute pairs to both sets. In a separate eight-person, four-person C3-to-C1 comparison on jtcxE69GiFV-13, A, B and C scored 72.914, 82.144 and 72.322. B retained six selected nodes, while A and C retained four; complete connection in this example used source-order review.

Among the seven of 84 images whose Q fell while prism IoU rose in the automatic twenty-person C2-to-C3 comparison, x8F5xyUWy9e-03 changed by $-$4.726 Q points, +2.450 IoU percentage points and +$0.836^\circ$ direction error. The complete outputs and component values explain the opposed quality summaries.

\paragraph{Cross-view structure and quality.}
Across the 128 directly related view pairs, same-person versus different-person mean gaps were 0.671 versus 0.752 effective turns, 0.377 versus 0.445 turn deficits and 0.363 versus 0.414 short-edge counts. Both comparisons use the same common-worker pool within a pair and equal pair weights. Across all pairs, simplification agreement was 88.22\% for the same person and 86.37\% for different people. Uniformly unsimplified responses left many pairs with little variation; the fifteen pairs with short edges in both references yielded the larger agreement difference reported in the main text.

Across four views of the illustrated room, A40 retained 10, 14, 16 and 16 effective turns, matching the reference counts; A41 retained 8, 4, 9 and twelve. Among eighteen common annotators in the two illustrated views, the door-side endpoint was matched in both views by five people, neither by six and one by seven; the vanity-side counts were five, eight and five. The counts agree at azimuth tolerances of 3, 6 and 9 degrees. The second manually confirmed room correspondence yielded inner-endpoint status changes for 19/33, 12/30 and 14/30 common annotators across three view pairs; its same-person endpoint agreement did not consistently exceed the different-person comparison.

Same-room mean-Q gaps averaged 5.051 points, compared with 8.931 for controls matched on building, difficulty, reference turns and geometrically hidden reference vertices. The comparison used common workers and 215 directional comparisons involving 113 same-room pairs in nine buildings. Gaps were smaller in eight buildings. Control pairs had unknown recorded relation and partially verified physical-room identity.

Within-pair Q rank correlation was positive on 105 of 128 pairs; forty of 57 people repeated relative-performance direction on more than half their pairs. Raw individual cross-view Q correlation was 0.401, peer-centered correlation 0.199 and correlation after adjustment for building-external personal history 0.013. The image-mean prediction uses both directions of 128 pairs, giving 256 comparisons; C7 defines its target-image and building aggregation and the separate individual-Q and structural prediction estimands.

The full-composition transfer analysis uses 324 panels from 127 directly related ordinary image pairs in thirteen buildings. Moments are averaged by direct pair, declared room group and building; aggregated covariance and variances yield the pooled correlation of $-$0.052. The specified replacement instead follows the same retained and exchanged people across nineteen view pairs in seven buildings. C14 reports automatic and assisted branches and their corresponding event supports.

\paragraph{Doorway shared-path comparison.}
On the fixed seventeen doorway images, screening excludes 145 of 584 submissions and retains 439. At ten people, the screening-minus-unscreened mean Q difference is +1.972 on the 3,411 of 4,352 paths scorable in both conditions. The main conditional-mean difference of +1.581 uses each condition's own scorable outputs. Both retain the same whole-space reference and underlying base orders. The complete 3--34-person doorway analysis uses decreasing image support above fourteen, whereas the main 3--14 curves retain all seventeen images.

\FloatBarrier

\subsection*{C20. Collection coverage and detailed study procedures}

\paragraph{Collection and analysis coverage.}

The primary personnel study uses 6,512 ordinary Manual responses from 60 people across 206 images in 22 buildings (130 easy, 39 moderate and 37 difficult images with 4,142, 1,231 and 1,139 responses, typically 31--35 responses per image). Table~\ref{tab:collection-coverage} places this set within the full collection. Offline teams recombine these independent human observations.



\begin{table}[!ht]
\centering\small
\caption{Response and image coverage by analysis set.}
\label{tab:collection-coverage}
\setlength{\tabcolsep}{4pt}\renewcommand{\arraystretch}{1.13}
\begin{tabular*}{\linewidth}{@{\extracolsep{\fill}}lrr@{}}
\hline
Analysis set & Responses & Images\\\hline
All submissions & 10,439 & ---\\
Primary-target quality & 8,154 & 225\\
\quad Ordinary Manual subset & 6,512 & 206\\
\quad Ordinary Semi subset & 960 & 37\\\hline
Alternative targets & 231 & 28\\
Doorway screening & 584 before; 439 after & 17\\
OOS diagnostics & 1,222 & 28\\\hline
\end{tabular*}
\end{table}

The table's ordinary-scene rows are subsets of the primary-target responses; the analysis sets overlap. Semi covers 59 people and 12 buildings. The alternative-target set comprises 31 pools. Initialization-error supports are specified with their results.



\paragraph{Historical profile inputs.}

\begin{table}[!ht]
\centering\small
\caption{Historical inputs and roles of annotator profiles.}
\label{tab:annotator-profiles}
\setlength{\tabcolsep}{3pt}\renewcommand{\arraystretch}{1.14}
\begin{tabular}{@{}>{\raggedright\arraybackslash}p{.22\linewidth}>{\raggedright\arraybackslash}p{.38\linewidth}>{\raggedright\arraybackslash}p{.33\linewidth}@{}}
\hline
Profile & Inputs & Question\\\hline
Structural expression & Turns, short edges, near-collinear points, reference matches & How do people draw?\\
Consensus response & Structural, boundary and scale changes in team outputs & How do members affect aggregation?\\
Joint profile & Structure, geometry, historical Q and time & How do expression, quality and labor relate?\\\hline
\end{tabular}
\end{table}

Profiles use known people's histories outside the target building. Target-building exclusion precedes feature construction, eligibility checks, scaling, dimension reduction and clustering. The structural profiles C1--C3 use structural-expression features; the consensus-response profiles R1--R2 summarize regularized member effects on group geometry. The four joint profiles J1--J4 use 13 features describing structure, boundary response, historical Q and recorded time. Geometric opportunity refers to reference turns, short edges and self-occlusion. Structural and joint profiles use domain-balanced features and Ward clustering. Within each training fold, the J labels follow increasing class-centre Q residual, with structural response and time resolving ties. Supplement C6--C9 specifies all features, class cuts and fitting procedures.



\paragraph{Feedback and cross-view procedures.}

Submission-time feedback records participants' reports of scene conditions and, for Semi, problems in the initialized proposal. Linked ordinary primary-target records cover 1,016 Manual submissions from 118 images and 456 Semi submissions from 34 images, with 24 people in each condition. Same-image contrasts compare people within each image and target; complementary same-person contrasts center outcomes within images and aggregate within person and building. Intervals resample buildings or people, respectively. Semi outcomes compare the submitted geometry with the exact initialization to identify retained errors and confirmed local corrections.

For cross-view structural repetition, we compare each annotator with themselves and with every other member of the same common-worker pool across directly related views. Simplification denotes fewer effective turns than the assigned reference. We report pair-equal comparisons and building-equal checks, including a subset whose two references both contain short edges. Local endpoint comparisons use established cross-view correspondences and one-to-one periodic azimuth matching at 3, 6 and 9 degrees. Structural prediction uses other directly related views, excludes the target view, and compares personal history with peers on the same source images. Errors are averaged over people within images, images within buildings, and then buildings.



\paragraph{Selection and stopping rules.}

Offline quality-and-time selection (Q/T) starts with the team ranked by building-external historical Q. It exchanges members for lower predicted labor while allowing at most a two-point decrease in mean historical relative Q. Each budget receives an independently selected team, which is then evaluated on its target image. The policy is compared with historical-Q selection and complements the structural-profile selection comparison (Supplement C2 and C8).

Source-history stopping uses the mean Q curve from all other available views of the same room. Along each shared member order, it chooses the earliest $k\geq3$ within two points of that source curve's peak. Target-view additions then continue until complete geometry is available. Comparators use a fixed seven-person request with the same completion rule, or stable node counts across three complete outputs. The replay evaluates both use of existing history and total labor including its acquisition. T sums selected members' recorded annotation durations in person-seconds; matched comparisons require compatible, positive finite time records. Supplement C8--C9 specifies the policies and accounting, and the main Discussion reports their results.



\paragraph{Outcome weighting and resampling.}

Q, its within-image tenth percentile (P10) and between-team standard deviation describe scorable layouts. Complete-and-evaluable rates use all attempted teams. Paired size changes require both outputs scorable; $G_S$ requires all three team scores, and $Q_C-Q_B$ requires B and C. Incomplete, failed and unscorable outputs keep their recorded status. Automatic and fixed-coordinate assisted-order supports are identified in each comparison.

Group-size and direct-replacement summaries generally weight images equally; selection results average contexts within images, images within buildings, and buildings equally. Group-size intervals use 2,000 bootstrap resamples of complete entry paths, preserving paired sizes and generation-batch sizes. These intervals describe Monte Carlo precision for the observed pools. Separate building resampling describes variation across the 22 sampled buildings while retaining the observed people. Supplement C17 gives the estimands, interpolation and resampling details. The resulting summaries distinguish variation across teams, computational precision and variation across buildings.

\end{document}

